\documentclass[11pt,a4paper]{article}

% ============================================================
%  Paquetes
% ============================================================
\usepackage[utf8]{inputenc}
\usepackage[T1]{fontenc}
\usepackage[spanish,es-tabla,es-nodecimaldot]{babel}
\usepackage{lmodern}
\usepackage[hmargin=2.5cm,vmargin=2.2cm,headsep=14pt]{geometry}
\usepackage{microtype}
\usepackage{amsmath}
\usepackage{booktabs}
\usepackage{tabularx}
\usepackage{float}

\usepackage[table]{xcolor}
\usepackage{enumitem}
\usepackage[font=small,labelfont=bf,labelsep=period]{caption}
\usepackage{fancyhdr}
\usepackage{tikz}
\usetikzlibrary{arrows.meta,positioning,calc,babel,decorations.pathreplacing}
\usepackage[hidelinks]{hyperref}

\setlist{itemsep=2pt, topsep=4pt}
\setlength{\parskip}{2pt}
\raggedbottom

% Encabezado y pie
\pagestyle{fancy}
\fancyhf{}
\fancyhead[L]{\small CS3014 -- Estructuras de Datos Avanzadas}
\fancyhead[R]{\small Propuesta de proyecto: Skip List vs.\ AVL Tree}
\fancyfoot[C]{\small\thepage}
\renewcommand{\headrulewidth}{0.4pt}

% Estilos de las figuras
\tikzset{
  >=Stealth,
  cell/.style={draw, minimum width=6mm, minimum height=4.2mm, inner sep=0pt},
  keycell/.style={cell, font=\small},
  tnode/.style={circle, draw, minimum size=6.5mm, inner sep=0pt, font=\small},
  stepbox/.style={draw, rounded corners=2pt, align=center, font=\small,
                  minimum height=10mm, text width=1.95cm, inner sep=3pt},
}
\newcommand{\on}{\cellcolor{gray!45}}

% ============================================================
\begin{document}

\thispagestyle{empty}
\begin{center}
  {\LARGE\bfseries Skip List frente a AVL Tree}\\[5pt]
  {\large una evaluación experimental sobre conjuntos ordenados dinámicos}\\[14pt]
  {\normalsize Renzo Josimar Felix Aponte \qquad Manyory Estefany Cueva Mendoza}\\[2pt]
  {\normalsize Jefersson Kevin Quicaña Erquinio}\\[6pt]
  {\small Universidad de Ingeniería y Tecnología (UTEC)}\\
  {\small CS3014 -- Estructuras de Datos Avanzadas \quad$\cdot$\quad Septiembre de 2026}
\end{center}
\vspace{2pt}
\noindent\rule{\textwidth}{0.4pt}

\begin{abstract}
\noindent
Proponemos implementar desde cero una Skip List y un AVL Tree y comparar
experimentalmente su tiempo de búsqueda, inserción y eliminación, con tamaños
de entrada entre $10^3$ y $10^6$ y datos en orden aleatorio, ordenado e
inverso. Ambas estructuras tienen costo $O(\log n)$ (esperado en la skip list,
de peor caso en el AVL); queremos ver cómo se comportan en la
práctica, sin partir de que una es mejor que la otra.
\end{abstract}
\noindent\rule{\textwidth}{0.4pt}

% ============================================================
\section{Introducción y motivación}

\paragraph{El problema.}
Muchos programas necesitan mantener un conjunto de claves ordenadas que cambia
con el tiempo: se insertan claves nuevas, se buscan claves existentes y se
eliminan otras. Esto aparece, por ejemplo, en el índice de una base de datos o en
el \texttt{std::set} de C++.

\paragraph{Las dos soluciones.}
La solución clásica son los árboles binarios de búsqueda balanceados, como el
AVL Tree, que garantizan $O(\log n)$ por operación en el peor caso a cambio
de aplicar rotaciones para mantenerse balanceados. En 1990, Pugh propuso la
\emph{skip list}~\cite{pugh1990} como alternativa: varias listas enlazadas
ordenadas, una encima de otra, donde la altura de cada nodo se decide al azar.
Así se obtiene $O(\log n)$ esperado sin rotaciones. Según el survey de
Xing, Vadrevu y Aref~\cite{vadrevu2024}, hoy la skip list se usa en sistemas
reales como LevelDB, RocksDB, Redis y la clase
\texttt{ConcurrentSkipListMap} de Java.

\paragraph{Por qué compararlas.}
Aunque dos estructuras tengan la misma complejidad asintótica, eso no dice cuál
es más rápida en la práctica. Las constantes ocultas, el número de punteros por nodo,
el uso de la memoria caché o el orden en que llegan los datos pueden cambiar el
tiempo real. Por eso queremos implementar ambas y medirlas bajo las mismas condiciones.


% ============================================================
\section{Conceptos y trabajo relacionado}\label{sec:conceptos}

\subsection{Skip List}

\paragraph{Idea.}
Una skip list está formada por varios niveles. El nivel 1 es una lista enlazada
ordenada con todas las claves. Cada nivel superior contiene solo una parte de
las claves del nivel de abajo y funciona como un ``carril rápido'' que permite
saltar varios nodos de una vez. Una cabecera tiene un puntero por nivel
(Figura~\ref{fig:skiplist}).

\begin{figure}[!htb]
\centering
\begin{tikzpicture}[x=1.25cm, y=4.2mm]
  \foreach \n/\x/\h/\k in {H/0/3/{}, a/1/1/3, b/2/2/5, c/3/1/8, d/4/3/12,
                           e/5/1/15, f/6/2/21, g/7/1/27, N/8/3/{}} {
    \ifx\k\empty\else \node[keycell] (\n-0) at (\x,0) {\k}; \fi
    \foreach \l in {1,...,\h} { \node[cell] (\n-\l) at (\x,\l) {}; }
  }
  \foreach \c in {H-3, d-3, d-2, d-1, e-0} {
    \fill[gray!25] (\c.south west) rectangle (\c.north east);
    \draw (\c.south west) rectangle (\c.north east); }
  \node[font=\small] at (e-0) {15};
  \node[font=\small] at (0,0) {cabecera};
  \node[draw, fill=white, minimum width=6mm, minimum height=12.6mm, font=\small] at (8,2) {NIL};
  \foreach \u/\v/\l in {H/a/1, a/b/1, b/c/1, c/d/1, d/e/1, e/f/1, f/g/1, g/N/1,
                        H/b/2, b/d/2, d/f/2, f/N/2, H/d/3, d/N/3} {
    \fill (\u-\l.center) circle (1.3pt);
    \draw[->] (\u-\l.center) -- (\v-\l.west |- \u-\l.center);
  }
  \foreach \l in {1,2,3} { \node[font=\scriptsize, anchor=west] at (8.35,\l) {nivel \l}; }
  \draw[->, line width=1.3pt] (H-3.center) -- (d-3.west);
  \draw[->, line width=1.3pt] (d-1.center) -- (e-1.west);
\end{tikzpicture}
\caption{Skip list con siete claves y $p = 1/2$. Las celdas sombreadas y los
punteros gruesos muestran la búsqueda de 15: se avanza por el nivel 3 hasta
12; como el siguiente nodo en los niveles 3 y 2 es mayor que 15, se baja hasta
el nivel 1, donde el siguiente nodo es 15.}
\label{fig:skiplist}
\end{figure}

\paragraph{Operaciones.}
\begin{itemize}
  \item \textbf{Búsqueda:} se empieza en el nivel más alto de la cabecera. En
  cada nivel se avanza mientras la siguiente clave sea menor que la buscada; si
  no, se baja un nivel. Se termina en el nivel 1.
  \item \textbf{Inserción:} se busca la posición guardando el último nodo
  visitado en cada nivel. Luego se decide la altura del nuevo nodo lanzando una
  ``moneda'': con probabilidad $p$ sube un nivel más, hasta un máximo
  \texttt{MaxLevel}. Finalmente se actualizan los punteros de los nodos
  guardados (Figura~\ref{fig:insert}).
  \item \textbf{Eliminación:} se busca la clave de la misma forma y se
  actualizan los punteros de sus predecesores en cada nivel donde aparece.
\end{itemize}

\begin{figure}[!htb]
\centering
\begin{tikzpicture}[x=1.12cm, y=4.2mm]
  \foreach \n/\x/\h/\k in {H/0/3/{}, a/1/1/3, b/2/2/5, c/3/1/8, d/4/3/12,
                           e/5/1/15, f/7/2/21, g/8/1/27, N/9/3/{}} {
    \ifx\k\empty\else \node[keycell] (\n-0) at (\x,0) {\k}; \fi
    \foreach \l in {1,...,\h} { \node[cell] (\n-\l) at (\x,\l) {}; }
  }
  % nodo nuevo (altura 2)
  \node[keycell, dashed, very thick] (x-0) at (6,0) {17};
  \node[cell, dashed, very thick] (x-1) at (6,1) {};
  \node[cell, dashed, very thick] (x-2) at (6,2) {};
  % predecesores cuyos punteros cambian
  \foreach \c in {d-2, e-1} {
    \fill[gray!25] (\c.south west) rectangle (\c.north east);
    \draw (\c.south west) rectangle (\c.north east); }
  \node[font=\small] at (0,0) {cabecera};
  \node[draw, fill=white, minimum width=6mm, minimum height=12.6mm, font=\small] at (9,2) {NIL};
  \foreach \u/\v/\l in {H/a/1, a/b/1, b/c/1, c/d/1, d/e/1, f/g/1, g/N/1,
                        H/b/2, b/d/2, f/N/2, H/d/3, d/N/3} {
    \fill (\u-\l.center) circle (1.3pt);
    \draw[->] (\u-\l.center) -- (\v-\l.west |- \u-\l.center);
  }
  % punteros nuevos
  \foreach \u/\v/\l in {d/x/2, e/x/1, x/f/1, x/f/2} {
    \fill (\u-\l.center) circle (1.3pt);
    \draw[->, line width=1.3pt] (\u-\l.center) -- (\v-\l.west |- \u-\l.center);
  }
  \node[font=\scriptsize, align=center, above=1mm of x-2] {nuevo nodo\\(altura 2)};
\end{tikzpicture}
\caption{Inserción de 17. La moneda sale ``sube'' una vez y ``no sube'' la
siguiente, así que el nodo tiene altura 2. Solo cambian los punteros de los
predecesores en esos niveles (celdas sombreadas: 12 en el nivel 2 y 15 en el
nivel 1); los punteros nuevos se muestran en trazo grueso.}
\label{fig:insert}
\end{figure}

\paragraph{Costo.}
Pugh~\cite{pugh1990} muestra que el número esperado de niveles es del orden de
$\log_{1/p} n$ y que el costo esperado de las tres operaciones es $O(\log n)$.
El número esperado de punteros por nodo es $1/(1-p)$, es decir, 2 cuando
$p = 1/2$, por lo que el espacio esperado es $O(n)$. El peor caso es $O(n)$,
pero tiene una probabilidad muy baja. Además, como la altura de cada nodo se
decide al azar, la forma de la estructura no depende del orden en que llegan
las claves.

\subsection{AVL Tree}

\paragraph{Idea.}
Un AVL Tree~\cite{avl1962} es un árbol binario de búsqueda con una regla extra:
para cada nodo, la diferencia de alturas entre su subárbol izquierdo y el
derecho es como máximo 1. Esa diferencia se llama \emph{factor de balance}.

\paragraph{Operaciones.}
La búsqueda es la misma que en cualquier árbol binario de búsqueda. La
inserción y la eliminación también empiezan igual, pero al terminar se
recorre el camino de vuelta hacia la raíz revisando el factor de balance. Si
algún nodo queda con factor $+2$ o $-2$, se aplica una rotación simple o doble
para recuperar el balance (Figura~\ref{fig:avl}).

\begin{figure}[!htb]
\centering
\begin{tikzpicture}[level distance=9mm, level 1/.style={sibling distance=15mm}]
  \begin{scope}
    \node[tnode] (r) {1}
      child[missing]
      child { node[tnode] {2}
        child[missing]
        child { node[tnode, fill=gray!25] {3} } };
    \node[font=\scriptsize, right=1mm of r] {$+2$};
    \node[font=\small, align=center] at (0.75,-2.75) {después de insertar\\1, 2, 3 (en orden)};
  \end{scope}
  \draw[->, thick] (2.4,-0.9) -- node[above, font=\small] {rotación} (4.1,-0.9);
  \begin{scope}[xshift=5.6cm, yshift=-0.45cm]
    \node[tnode] {2}
      child { node[tnode] {1} }
      child { node[tnode, fill=gray!25] {3} };
    \node[font=\small, align=center] at (0,-2.3) {árbol AVL\\balanceado};
  \end{scope}
\end{tikzpicture}
\caption{Al insertar claves en orden creciente, el AVL pierde el balance en la
raíz (factor $+2$) y lo recupera con una rotación simple a la izquierda.}
\label{fig:avl}
\end{figure}

\paragraph{Costo.}
Gracias a la regla de balance, la altura del árbol es como máximo
aproximadamente $1.44 \log_2(n+2)$~\cite{knuth1998}, así que búsqueda,
inserción y eliminación cuestan $O(\log n)$ en el \emph{peor caso}. Cada nodo
guarda la clave, dos punteros y su altura, por lo que el espacio es $O(n)$. A
diferencia de la skip list, la forma del árbol sí depende del orden en que
llegan las claves, aunque las rotaciones siempre lo mantienen balanceado.

\subsection{Comparación teórica}

La Tabla~\ref{tab:teoria} junta lo que ya se conoce en la literatura. Las dos
estructuras tienen el mismo orden de complejidad, pero una da una garantía
esperada y la otra de peor caso, y mantienen el balance de forma distinta.
La teoría no dice cuál será más rápida en una computadora real, y eso es lo
que queremos medir.

\begin{table}[!htb]
\centering
\begin{tabular}{lcc}
\toprule
\textbf{Aspecto} & \textbf{Skip List} & \textbf{AVL Tree} \\
\midrule
Búsqueda, inserción, eliminación & $O(\log n)$ esperado & $O(\log n)$ peor caso \\
Peor caso de una operación       & $O(n)$ (poco probable) & $O(\log n)$ \\
Espacio                          & $O(n)$ esperado & $O(n)$ \\
Cómo mantiene el balance         & azar (altura aleatoria) & rotaciones \\
¿Depende del orden de los datos? & no & sí (cambian las rotaciones) \\
\bottomrule
\end{tabular}
\caption{Comparación teórica según la literatura~\cite{pugh1990,knuth1998}.}
\label{tab:teoria}
\end{table}

\subsection{Trabajo relacionado}

El trabajo base es el artículo de Pugh~\cite{pugh1990}, que introduce la skip
list, analiza su costo esperado y discute la elección de $p$: recomienda
$p = 1/2$ en general y $p = 1/4$ si se quiere menos variabilidad y menos
memoria. Pugh también reporta comparaciones de tiempo frente a árboles
balanceados, pero con el hardware y los compiladores de 1990, por lo que no
necesariamente representan lo que pasa en una computadora actual.

Como referencia complementaria usamos el survey de Xing, Vadrevu y
Aref~\cite{vadrevu2024}, que revisa muchas variantes de la skip list
(deterministas, concurrentes, orientadas a caché y a memoria externa) y su uso
en bases de datos. Entre las deterministas está la 1-2 skip list de Munro,
Papadakis y Sedgewick~\cite{munro1992}, equivalente a un árbol 2-3. En este
proyecto no implementaremos variantes: nos quedamos con la versión básica de
Pugh.

% ============================================================
\section{Planteamiento del proyecto}\label{sec:plan}

\subsection{Pregunta de investigación}

\begin{quote}
\itshape
¿Cómo se comporta una Skip List frente a un AVL Tree en operaciones de
búsqueda, inserción y eliminación sobre conjuntos ordenados dinámicos?
\end{quote}

Para poder responderla con experimentos, la dividimos en dos preguntas más
concretas:
\begin{itemize}
  \item[\textbf{P1.}] ¿Cómo cambia el tiempo de cada operación cuando crece el
  tamaño de la entrada $n$?
  \item[\textbf{P2.}] ¿Influye el orden de los datos (aleatorio, ordenado o
  inverso) en el tiempo de cada estructura?
\end{itemize}

\subsection{Objetivos}

\paragraph{Objetivo general.}
Comparar experimentalmente el tiempo de ejecución de una Skip List y de un AVL
Tree, ambos implementados desde cero, en búsqueda, inserción y eliminación,
bajo distintos tamaños de entrada y órdenes de los datos.

\paragraph{Objetivos específicos.}
\begin{itemize}
  \item[\textbf{O1.}] Implementar una Skip List según el diseño de Pugh~\cite{pugh1990}.
  \item[\textbf{O2.}] Implementar un AVL Tree con inserción y eliminación balanceadas.
  \item[\textbf{O3.}] Verificar la correctitud de ambas implementaciones con
  pruebas automáticas.
  \item[\textbf{O4.}] Construir un generador de datos y un programa de
  medición que apliquen las mismas operaciones a ambas estructuras.
  \item[\textbf{O5.}] Medir, graficar y analizar los tiempos junto con la
  complejidad teórica de cada estructura.
\end{itemize}

% ============================================================
\section{Metodología experimental}\label{sec:metodo}

La Tabla~\ref{tab:diseno} resume el diseño experimental; las subsecciones
siguientes explican cada parte.

\begin{table}[!htb]
\centering
\renewcommand{\arraystretch}{1.2}
\begin{tabularx}{\textwidth}{l X}
\toprule
\textbf{Elemento} & \textbf{Decisión} \\
\midrule
Estructuras   & Skip List ($p=1/2$, \texttt{MaxLevel}$\,=20$) y AVL Tree, ambas en C++17 desde cero \\
Claves        & Enteros de 64 bits, sin duplicados \\
Operaciones   & Inserción, búsqueda (50\,\% claves presentes, 50\,\% ausentes) y eliminación \\
Tamaños       & $n \in \{10^3,\ 10^4,\ 10^5,\ 10^6\}$ \\
Orden de datos& Aleatorio, ordenado (creciente) e inverso (decreciente) \\
Escenarios    & $4 \text{ tamaños} \times 3 \text{ órdenes} = 12$ por estructura \\
Repeticiones  & 10 por escenario, con semillas distintas \\
Métrica       & Tiempo promedio por operación (ns) y desviación estándar \\
Entorno       & Misma máquina y mismas opciones de compilación (\texttt{-O2}) \\
\bottomrule
\end{tabularx}
\caption{Resumen del diseño experimental.}
\label{tab:diseno}
\end{table}

\subsection{Implementación y verificación}

Ambas estructuras tendrán la misma interfaz: \texttt{insert(k)},
\texttt{search(k)} y \texttt{remove(k)}. No usaremos librerías externas para
las estructuras, para controlar las condiciones del experimento. Elegimos
$\texttt{MaxLevel} = 20$ porque $\log_2 10^6 \approx 20$ alcanza para el
tamaño máximo. Si el tiempo lo permite, repetiremos una parte de los
experimentos con $p = 1/4$.

Antes de medir tiempos verificaremos que ambas implementaciones sean correctas:
ejecutaremos secuencias aleatorias de operaciones y compararemos los resultados
con \texttt{std::set}, que se usará solo como referencia en las pruebas y no en
las mediciones. También revisaremos que el AVL cumpla su condición de balance y
que cada nivel de la Skip List esté ordenado.

\subsection{Procedimiento de medición}

La Figura~\ref{fig:pipeline} muestra el procedimiento que se aplicará a cada
combinación de tamaño y orden de datos.

\begin{figure}[!htb]
\centering
\begin{tikzpicture}[node distance=4.5mm]
  \node[stepbox] (gen) {Generar $n$ claves\\(semilla fija)};
  \node[stepbox, right=of gen] (ins) {\textbf{Insertar}\\$n$ claves};
  \node[stepbox, right=of ins] (bus) {\textbf{Buscar}\\$n$ consultas};
  \node[stepbox, right=of bus] (eli) {\textbf{Eliminar}\\$n$ claves};
  \node[stepbox, right=of eli] (rep) {Repetir\\10 veces};
  \node[stepbox, right=of rep] (est) {Promedio y\\gráficos};
  \foreach \u/\v in {gen/ins, ins/bus, bus/eli, eli/rep, rep/est} \draw[->] (\u) -- (\v);
  \draw[->, dashed] (rep.south) -- ++(0,-4mm) -| node[pos=0.25, below, font=\scriptsize] {nuevas semillas} (gen.south);
  \draw[decorate, decoration={brace, amplitude=4pt}]
    ($(ins.north west)+(0,2mm)$) -- ($(eli.north east)+(0,2mm)$)
    node[midway, above=4pt, font=\scriptsize] {se mide el tiempo de cada operación};
\end{tikzpicture}
\caption{Procedimiento experimental para cada tamaño $n$ y orden de datos.
Ambas estructuras reciben exactamente las mismas claves y consultas.}
\label{fig:pipeline}
\end{figure}

En cada repetición, la \textbf{inserción} parte de una estructura vacía y
recibe las $n$ claves en el orden correspondiente; la \textbf{búsqueda}
ejecuta $n$ consultas en orden aleatorio sobre la estructura ya construida; y
la \textbf{eliminación} borra las $n$ claves en el mismo orden de la
inserción.

Para reducir el ruido de las mediciones, se descartará una primera ejecución
de calentamiento, no se imprimirá nada mientras se mide y el tiempo se tomará
con \texttt{std::chrono::steady\_clock} sobre cada lote completo de
operaciones. Opcionalmente, registraremos también la altura final del AVL y el
nivel máximo alcanzado por la Skip List, si ayudan a explicar los resultados.

\subsection{Cómo se presentarán los resultados}

Cada gráfico corresponde a una de las preguntas de la Sección~\ref{sec:plan} (Tabla~\ref{tab:graficos}). Los valores se obtendrán
únicamente durante la experimentación.

\begin{table}[!htb]
\centering
\renewcommand{\arraystretch}{1.2}
\begin{tabularx}{\textwidth}{l l l X}
\toprule
\textbf{Gráfico} & \textbf{Eje $x$} & \textbf{Eje $y$} & \textbf{Qué responde} \\
\midrule
G1: Búsqueda    & $n$ (escala log) & ns por operación & P1 para la búsqueda, Skip List vs.\ AVL \\
G2: Inserción   & $n$ (escala log) & ns por operación & P1 para la inserción, Skip List vs.\ AVL \\
G3: Eliminación & $n$ (escala log) & ns por operación & P1 para la eliminación, Skip List vs.\ AVL \\
G4: Orden       & orden de datos   & ns por operación & P2: efecto del orden con $n = 10^6$ \\
\bottomrule
\end{tabularx}
\caption{Gráficos planeados y la pregunta que responde cada uno.}
\label{tab:graficos}
\end{table}

% ============================================================
\section{Plan de trabajo}\label{sec:cronograma}

El proyecto va de la semana 9 a la 15 del ciclo. En la semana 9 se entrega
esta propuesta, las semanas 10 a 14 son las cinco semanas de desarrollo y la
semana 15 corresponde a la exposición final. La Tabla~\ref{tab:cronograma}
muestra cómo repartimos el trabajo y el resultado que esperamos tener al final
de cada semana.

\begin{table}[!htb]
\centering
\renewcommand{\arraystretch}{1.25}
\small
\setlength{\tabcolsep}{4pt}
\begin{tabular}{>{\raggedright\arraybackslash}p{4.7cm}*{7}{>{\centering\arraybackslash}p{0.62cm}}>{\raggedright\arraybackslash}p{4.4cm}}
\toprule
\textbf{Actividad} & \textbf{S9} & \textbf{S10} & \textbf{S11} & \textbf{S12} & \textbf{S13} & \textbf{S14} & \textbf{S15} & \textbf{Resultado} \\
\midrule
Propuesta del proyecto            & \on & & & & & & & Propuesta entregada \\
Revisión bibliográfica y diseño   & & \on & & & & & & Interfaces y plan de pruebas \\
Implementación del código         & & & \on & & & & & Código de ambas estructuras \\
Pruebas y datos de prueba         & & & & \on & & & & Pruebas y generador listos \\
Experimentos y medición           & & & & & \on & & & Tiempos de los 12 escenarios \\
Análisis, gráficos y redacción    & & & & & & \on & & Gráficos G1--G4 e informe \\
Exposición final                  & & & & & & & \on & Presentación \\
\bottomrule
\end{tabular}
\caption{Cronograma tentativo por semana del ciclo (S = semana).}
\label{tab:cronograma}
\end{table}

% ============================================================
\section{Entregable y resultados esperados}\label{sec:resultados}

\paragraph{Entregables.}
Al finalizar el proyecto esperamos entregar:
\begin{itemize}
  \item El código fuente de la Skip List y del AVL Tree, con sus pruebas.
  \item Los scripts para generar los datos, ejecutar los experimentos y
  producir los gráficos.
  \item Los gráficos G1--G4 y las tablas con los tiempos medidos.
  \item Un informe final que compare los resultados con la complejidad teórica.
\end{itemize}

No sabemos de antemano qué estructura será más rápida, y no queremos
suponerlo; los resultados de P1 y P2 lo dirán.

\paragraph{Posibles dificultades.}
Las mediciones de tiempo pueden variar por el sistema operativo o la memoria
caché, por eso usaremos varias repeticiones y la misma máquina para todo.
También creemos que la eliminación en el AVL es la parte más fácil de
implementar mal, así que será lo primero que probemos contra
\texttt{std::set}. Por último, los experimentos con $n = 10^6$ pueden
tardar, así que automatizaremos las corridas con scripts desde la semana 12.

% ============================================================
\begin{thebibliography}{9}
\small
\setlength{\itemsep}{1pt}

\bibitem{pugh1990}
W.~Pugh.
\newblock Skip lists: A probabilistic alternative to balanced trees.
\newblock \emph{Communications of the ACM}, 33(6):668--676, 1990.
\newblock \url{https://doi.org/10.1145/78973.78977}

\bibitem{vadrevu2024}
L.~Xing, V.~S.~P.~K.~Vadrevu y W.~G.~Aref.
\newblock The ubiquitous skiplist: A survey of what cannot be skipped about the
skiplist and its applications in data systems.
\newblock \emph{ACM Computing Surveys}, 57(11):1--37, 2025.
\newblock \url{https://doi.org/10.1145/3736754}

\bibitem{avl1962}
G.~M.~Adel'son-Vel'skii y E.~M.~Landis.
\newblock An algorithm for the organization of information.
\newblock \emph{Soviet Mathematics Doklady}, 3:1259--1263, 1962.

\bibitem{knuth1998}
D.~E.~Knuth.
\newblock \emph{The Art of Computer Programming, Vol.~3: Sorting and
Searching}, 2.\textsuperscript{a} ed.
\newblock Addison-Wesley, 1998.

\bibitem{munro1992}
J.~I.~Munro, T.~Papadakis y R.~Sedgewick.
\newblock Deterministic skip lists.
\newblock En \emph{Proceedings of the Third Annual ACM-SIAM Symposium on
Discrete Algorithms (SODA)}, pp.~367--375, 1992.

\end{thebibliography}

\end{document}
